% Source: https://artofproblemsolving.com/community/c3h1611099_national_mathcounts_problem
% Stable identity: aops-topic-1611099-q01
\textbf{A tetrahedron with edge $\sqrt2$ fits.} Put the cube at $[0,1]^3$ and take four alternate vertices $A=(0,0,0)$, $B=(1,1,0)$, $C=(1,0,1)$, $D=(0,1,1)$. Each pair of them differ in exactly two coordinates, so every edge is a face diagonal of length $\sqrt2$, and $ABCD$ is a regular tetrahedron inside the cube.

\begin{asy}
usepackage("amsmath");
size(5cm);
defaultpen(fontsize(10pt));
pair O=(0,0),X=(3,0),Y=(1.1,0.75),Z=(0,3);
pair P(real a,real b,real c){return a*X+b*Y+c*Z;}
pair A=P(0,0,0),B=P(1,1,0),C=P(1,0,1),D=P(0,1,1);
draw(P(0,0,0)--P(1,0,0)--P(1,0,1)--P(0,0,1)--cycle);
draw(P(1,0,0)--P(1,1,0)--P(1,1,1)--P(1,0,1));
draw(P(0,0,1)--P(0,1,1)--P(1,1,1));
draw(P(0,0,0)--P(0,1,0)--P(1,1,0),dashed);draw(P(0,1,0)--P(0,1,1),dashed);
draw(A--C--D--A,blue+linewidth(0.9));draw(C--B,blue+linewidth(0.9));
draw(A--B,blue+dashed);draw(B--D,blue+dashed);
dot(A);dot(B);dot(C);dot(D);
label("$A$",A,SW);label("$B$",B,E);label("$C$",C,(1,-0.6));label("$D$",D,N);
label("$1$",(P(0,0,0)+P(1,0,0))/2,S);
label("\mbox{}",(min(currentpicture,user=true)+max(currentpicture,user=true))/2,black); // keeps the text after the figure black when LaTeX typesets the labels inline
\end{asy}

\textbf{No larger one fits.} Let a regular tetrahedron with edge $s$ lie in the cube, with vertices $P_1,\dots,P_4$. The sum of the squares of its six edges is
\[
6s^2=\sum_{i<j}|P_iP_j|^2=\sum_{\text{coordinates}}\ \sum_{i<j}\left(t_i-t_j\right)^2,
\]
where for each of the three coordinates $t_1,\dots,t_4\in[0,1]$ are the values of that coordinate at the four vertices. For four numbers in $[0,1]$,
\[
\sum_{i<j}\left(t_i-t_j\right)^2=4\sum t_i^2-\left(\sum t_i\right)^2 .
\]
As a function of any one $t_i$ (the others fixed) this is a quadratic with positive leading coefficient $4-1=3$, so it is largest at an endpoint $t_i\in\{0,1\}$. If $k$ of the four numbers are $1$ and the rest $0$, the sum is $k(4-k)\le4$. So each coordinate contributes at most $4$, giving $6s^2\le12$ and $s\le\sqrt2$.

\textbf{Surface area.} The largest tetrahedron has edge $\sqrt2$, so each face is an equilateral triangle of area $\dfrac{\sqrt3}{4}\left(\sqrt2\right)^2=\dfrac{\sqrt3}{2}$, and the total surface area is
\[
4\cdot\dfrac{\sqrt3}{2}=\boxed{2\sqrt3\text{ cm}^2}.
\]
